\documentclass[12pt, titlepage]{article}

\usepackage{amsmath, mathtools}

\usepackage[round]{natbib}
\usepackage{amsfonts}
\usepackage{amssymb}
\usepackage{graphicx}
\usepackage{colortbl}
\usepackage{xr}
\usepackage{hyperref}
\usepackage{longtable}
\usepackage{xfrac}
\usepackage{tabularx}
\usepackage{float}
\usepackage{siunitx}
\usepackage{booktabs}
\usepackage{multirow}
\usepackage[section]{placeins}
\usepackage{caption}
\usepackage{fullpage}

\hypersetup{
bookmarks=true,     % show bookmarks bar?
colorlinks=true,       % false: boxed links; true: colored links
linkcolor=red,          % color of internal links (change box color with linkbordercolor)
citecolor=blue,      % color of links to bibliography
filecolor=magenta,  % color of file links
urlcolor=cyan          % color of external links
}

\usepackage{array}

% Requirements are referenced by stable SRS identifiers and source links.

\input{../../Comments.text}
\input{../../Common.text}

\begin{document}

\title{Module Interface Specification for \progname{}}

\author{\authname}

\date{\today}

\maketitle

\pagenumbering{roman}

\section{Revision History}

\begin{tabularx}{\textwidth}{p{3cm}p{2cm}X}
\toprule {\bf Date} & {\bf Version} & {\bf Notes}\\
\midrule
November 13, 2026 & -1 & Initial module inventory and first-half interface syntax.\\
\bottomrule
\end{tabularx}

~\newpage

\section{Symbols, Abbreviations and Acronyms}

The authoritative functional and nonfunctional identifiers (FR1--FR16,
IN1--IN12, etc.) are defined in the
\href{https://github.com/rhea333/LingoLyrics/blob/main/docs/SRS-Volere/SRS.tex}
{LingoLyrics Volere SRS}. M1--M8 refer to the module numbers in the
Module Guide; MIS means Module Interface Specification, and ADT means
abstract data type. Terms such as \emph{song map}, \emph{line},
\emph{annotation}, and \emph{target} follow the definitions in the SRS.

\newpage

\tableofcontents

\newpage

\pagenumbering{arabic}

\section{Introduction}

This Module Interface Specification describes the externally observable
contracts of LingoLyrics's software modules. The mobile learner
application opens and plays locally stored lessons; the browser editor
creates validated song-map packages. Both depend on shared map and
validation services. The architecture and each module's hidden design
decisions are documented in the
\href{https://github.com/rhea333/LingoLyrics/blob/docs/mg/docs/Design/SoftArchitecture/MG.tex}
{Module Guide}. Source code and project documents are in the
\href{https://github.com/rhea333/LingoLyrics}{LingoLyrics repository}.

An MIS defines the behaviour, state and environmental interactions that
a consumer may depend on; it does not prescribe UI widget structure,
runtime library calls, or a particular serialization language.
The interface specifications for M1--M4 are provided in this revision.
M5--M8 require companion specifications before the team's Revision~$-1$
submission can be considered complete.

\section{Notation}

This specification uses abstract types and simple mathematical contracts
instead of application-specific programming syntax. A function
$f : A \rightarrow B$ accepts a value of type $A$ and yields a value
of type $B$. $\mathbb{R}_{\geq 0}$ is the set of nonnegative real
numbers, $\mathbb{B}=\{\mathit{true},\mathit{false}\}$, and
$\operatorname{Seq}(T)$, $\operatorname{Set}(T)$, and
$\operatorname{Option}(T)$ mean a sequence, set, and optional value of
type $T$. $\bot$ denotes no available value. Time measurements
are in seconds; $t$ is a timestamp and $d$ is audio duration.

Types from the shared format module M6 include
$\mathrm{SongMap}$, $\mathrm{LineRef}$, $\mathrm{TargetRef}$,
$\mathrm{VocabularyItem}$, and $\mathrm{LessonID}$.
A $\mathrm{CandidateMap}$ is an untrusted, potentially invalid input
to the validator; a $\mathrm{ValidatedMap}$ is a candidate accepted
by all relevant song-map validity rules (SRS V1--V8).
Other interface-level types used below are
$\mathrm{AudioRef}$ (a reference to an audio resource),
$\mathrm{FileRef}$ (an opaque file location),
$\mathrm{Bytes}$ (an opaque file payload),
$\mathrm{Annotation}$ (lesson context and translation),
$\mathrm{Edit}$ (one undoable change),
$\mathrm{RuleID}$ (a validation rule identifier), and
$\mathrm{ValidationResult}$ (a Boolean verdict with failed rule IDs).

The \emph{Uses} subsection lists only the modules whose public
services are required to define this interface, not every
implementation dependency. \emph{State variables} belong to exactly
one module. \emph{Environment variables} represent devices, files,
displays, and external services not controlled by the module.
For access program semantics, $s' := v$ denotes a new state after
the operation, $out := v$ its returned result, and $exc := E$ an
exception or failure. Unless specified otherwise, an operation does
not mutate its input values.

\section{Module Decomposition} \label{SecMISModules}

The modules and numbers are taken from the
\href{https://github.com/rhea333/LingoLyrics/blob/docs/mg/docs/Design/SoftArchitecture/MG.tex}
{Module Guide}. The table is a classification by hidden decision,
not a call graph.

\begin{table}[H]
\centering
\begin{tabularx}{\textwidth}{@{}p{0.27\textwidth}lX@{}}
\toprule
\textbf{Class} & \textbf{ID} & \textbf{Module}\\
\midrule
Hardware-hiding & M1 & Platform and Media Adapter\\
Behaviour-hiding & M2 & Learner Lesson Interface\\
 & M3 & Song Map Authoring Interface\\
 & M4 & Song Map Validation Module\\
 & M5 & Vocabulary Handoff Module\\
Software decision & M6 & Shared Song Map Format Module\\
 & M7 & Playback Synchronization Module\\
 & M8 & Language Text Processing Module\\
\bottomrule
\end{tabularx}
\caption{LingoLyrics module inventory matching the Module Guide}
\label{TblMH}
\end{table}

\clearpage
\section{MIS of M1: Platform and Media Adapter} \label{MISM1}

\subsection{Module}
Platform and Media Adapter (hardware-hiding abstract object).
A single adapter instance represents the active platform playback
session and access to local files. The adapter hides OS-specific
audio and storage operations from consumer modules.

\subsection{Uses}
None.

\subsection{Syntax}
\subsubsection{Exported Constants}
None. Playback tolerances are defined as requirement parameters,
not hard-coded constants of the adapter.

\subsubsection{Exported Access Programs}
\begin{center}
\small
\begin{tabularx}{\textwidth}{@{}p{0.20\textwidth}p{0.25\textwidth}p{0.23\textwidth}X@{}}
\toprule
\textbf{Name} & \textbf{In} & \textbf{Out} & \textbf{Exceptions}\\
\midrule
\texttt{loadAudio} & $\mathrm{AudioRef}$ & -- &
ResourceMissing, UnsupportedAudio\\
\texttt{play} & -- & -- & NoAudio\\
\texttt{pause} & -- & -- & NoAudio\\
\texttt{seek} & $t:\mathbb{R}_{\geq 0}$ & -- & NoAudio, OutOfRange\\
\texttt{getPosition} & -- & $\mathbb{R}_{\geq 0}$ & NoAudio\\
\texttt{getDuration} & -- & $\mathbb{R}_{\geq 0}$ & NoAudio\\
\texttt{readFile} & $\mathrm{FileRef}$ & $\mathrm{Bytes}$ &
ResourceMissing, AccessDenied\\
\texttt{writeFile} & $\mathrm{FileRef},\mathrm{Bytes}$ & -- &
AccessDenied, NoSpace\\
\bottomrule
\end{tabularx}
\normalsize
\end{center}

\subsection{Semantics}
\subsubsection{State Variables}
$activeAudio : \operatorname{Option}(\mathrm{AudioRef})$.
Initially $activeAudio=\bot$. This is the sole logical owner of the
currently loaded audio resource; UI modules do not copy playback
engine state.

\subsubsection{Environment Variables}
$devicePosition:\mathbb{R}_{\geq0}$ and
$deviceDuration:\mathbb{R}_{\geq0}$ are reported by the external
audio device; $deviceMode\in\{\mathrm{Idle},\mathrm{Ready},
\mathrm{Playing},\mathrm{Paused},\mathrm{Ended}\}$ describes its
transport state. $fileStore$ denotes the external platform filesystem.
These are observations or effects outside the module, not private
software state variables.

\subsubsection{Assumptions and Invariants}
A loaded audio resource has a finite positive duration.
When $activeAudio\neq\bot$, the reported position lies in
$[0,deviceDuration]$, subject to the device's reporting precision.
Only M1 may issue commands to its active audio device.
The caller is responsible for the user-action rules in HS1;
loading a resource alone never starts playback.

\subsubsection{Access Routine Semantics}
\begin{description}
\item[\texttt{loadAudio}(a):] On successful decode, $activeAudio':=a$,
$devicePosition':=0$, and $deviceMode':=\mathrm{Ready}$.
Otherwise $exc:=\mathrm{ResourceMissing}$ or
$\mathrm{UnsupportedAudio}$; the old session remains unchanged.
\item[\texttt{play}():] If no audio is loaded, $exc:=\mathrm{NoAudio}$;
otherwise $deviceMode':=\mathrm{Playing}$.
\item[\texttt{pause}():] If no audio is loaded, $exc:=\mathrm{NoAudio}$;
otherwise $deviceMode':=\mathrm{Paused}$.
\item[\texttt{seek}(t):] If there is no active audio, raise
NoAudio; if $t>deviceDuration$, raise OutOfRange; otherwise set the
external audio position to $t$.
\item[\texttt{getPosition}(), \texttt{getDuration}():]
After a successful load, $out:=devicePosition$ and
$out:=deviceDuration$, respectively. NoAudio is raised otherwise.
\item[\texttt{readFile}(f):] $out:=fileStore[f]$ when readable;
raise ResourceMissing or AccessDenied otherwise.
\item[\texttt{writeFile}(f,b):] On success, $fileStore'[f]:=b$.
Raise AccessDenied or NoSpace on failure; do not report success
before the requested write is complete.
\end{description}

\clearpage
\section{MIS of M2: Learner Lesson Interface} \label{MISM2}

\subsection{Module}
Learner Lesson Interface (behaviour-hiding abstract object). It owns
lesson selection, presentation preferences, and learner interactions;
it does not implement the audio clock, validator, tokenizer, or
Anki transport.

\subsection{Uses}
M1 Platform and Media Adapter, M4 Song Map Validation,
M5 Vocabulary Handoff, M6 Shared Song Map Format, and
M7 Playback Synchronization.

\subsection{Syntax}
\subsubsection{Exported Constants}
None. UI thresholds and lesson limits are taken from the SRS.

\subsubsection{Exported Access Programs}
\begin{center}
\small
\begin{tabularx}{\textwidth}{@{}p{0.21\textwidth}p{0.24\textwidth}p{0.22\textwidth}X@{}}
\toprule
\textbf{Name} & \textbf{In} & \textbf{Out} & \textbf{Exceptions}\\
\midrule
\texttt{importLesson} & $\mathrm{FileRef}$ &
$\mathrm{LessonID}$ & InvalidMap, ResourceMissing\\
\texttt{openLesson} & $\mathrm{LessonID}$ & -- & UnknownLesson\\
\texttt{play}, \texttt{pause} & -- & -- & NoLesson\\
\texttt{seek} & $t:\mathbb{R}_{\geq0}$ & -- & NoLesson, OutOfRange\\
\texttt{selectTarget} & $\mathrm{TargetRef}$ &
$\operatorname{Option}(\mathrm{Annotation})$ & UnknownTarget\\
\texttt{setLayers} & $\mathbb{B}\times\mathbb{B}$ & -- & --\\
\texttt{sendSelected} & -- & $\mathrm{SendStatus}$ &
NoSelection\\
\bottomrule
\end{tabularx}
\normalsize
\end{center}

\subsection{Semantics}
\subsubsection{State Variables}
$library:\operatorname{Seq}(\mathrm{LessonID})$ is the locally
persisted lesson index; $opened:
\operatorname{Option}(\mathrm{ValidatedMap})$ is the active
lesson; $selected:\operatorname{Option}(\mathrm{TargetRef})$
is the selected word or line; and
$showTransliteration,showTranslation:\mathbb{B}$ are remembered
presentation preferences. Initially $opened=selected=\bot$.
The library and preferences are restored from local storage.

\subsubsection{Environment Variables}
$learnerScreen$ is the displayed lesson and its controls;
$learnerInput$ is the stream of taps, swipes, and accessibility
actions. File storage and the audio player belong to M1 and are
observed through that module rather than duplicated as M2 state.

\subsubsection{Assumptions and Invariants}
$selected\neq\bot\Rightarrow opened\neq\bot$.
Every $opened$ lesson has passed M4 validation and its major
version is supported (FR1, IN3). Displayed segmentation is read
from the validated map and is never recomputed by M2 (IN4).
M2 must not display a successful Anki transfer before the M5
service confirms it (IN8).

\subsubsection{Access Routine Semantics}
\begin{description}
\item[\texttt{importLesson}(f):] Read the candidate from M1 and
validate it using M4. On success, append its identifier to
$library$ if not already present and persist the updated index;
$out:=\mathrm{LessonID}$. An invalid candidate raises InvalidMap
and leaves $library$ unchanged.
\item[\texttt{openLesson}(id):] If $id$ is in the library, load
the validated map and set $opened'$ to it and $selected':=\bot$;
otherwise raise UnknownLesson. The screen presents that lesson.
\item[\texttt{play}(), \texttt{pause}(), \texttt{seek}(t):]
When $opened=\bot$, raise NoLesson. Otherwise delegate to M1
and use the current M1 position with M7 to identify the active
lyric. Playback starts only on a learner action (HS1).
\item[\texttt{selectTarget}(r):] If $r$ is not in the active
lesson, raise UnknownTarget. Otherwise $selected':=r$ and
$out$ is the stored optional annotation; playback continues.
\item[\texttt{setLayers}(b_1,b_2):]
$showTransliteration':=b_1$ and $showTranslation':=b_2$;
persist these preferences on the device.
\item[\texttt{sendSelected}():] Require a selected word. Pass
its term, correct word translation, and source-line context to
M5; return its confirmed, queued, or failed result without
claiming success on a mere attempted send.
\end{description}

\bibliographystyle {plainnat}
\bibliography {../../../refs/References}

\newpage

\newpage{}

\section{Appendix --- Generative AI Use Disclosure}

\wss{ChatGPT (version 5) was used to brainstorm ideas for the template for this disclose section.}

\subsection{AI Tools Used}

\wss{Give the name of the AI tool(s) used and the version.}

\subsection{How and Where Used}

\wss{Explain what the tool(s) were used for.  Examples include: brainstorming,
explaining a technical concept, reviewing the document, generating or modifying
text, generating or modifying diagrams, identifying errors or omissions or
inconsistencies.}

\wss{For each of the uses, identify which parts of the document were affected.}

\wss{You do not need to report every interaction, but you should disclose the 
uses that materially contributed to your work}

\subsection{Verification of AI-Generated Content}

\wss{\begin{itemize}
    \item Describe how the team checked AI-generated or AI-assisted material
    before including it in the document.
    \item In particular, verify factual claims, technical assertions,
    requirements, calculations, references, and other information for which an
    AI system may produce incorrect or fabricated results.
    \item \textbf{The team is responsible for the accuracy, completeness,
consistency, and appropriateness of the submitted document, regardless of
whether material was generated by a human or an AI system.}
\end{itemize}}

\newpage{}

\section{Appendix --- Reflection}

\wss{Not required for CAS 741 projects}

The information in this section will be used to evaluate the team members on the
graduate attribute of Problem Analysis and Design.

\input{../../Reflection.text}

\begin{enumerate}
  \item What went well while writing this deliverable? 
  \item What pain points did you experience during this deliverable, and how
    did you resolve them?
  \item Which of your design decisions stemmed from speaking to your client(s)
  or a proxy (e.g. your peers, stakeholders, potential users)? For those that
  were not, why, and where did they come from?
  \item While creating the design doc, what parts of your other documents (e.g.
  requirements, hazard analysis, etc), it any, needed to be changed, and why?
  \item What are the limitations of your solution?  Put another way, given
  unlimited resources, what could you do to make the project better? (LO\_ProbSolutions)
  \item Give a brief overview of other design solutions you considered.  What
  are the benefits and tradeoffs of those other designs compared with the chosen
  design?  From all the potential options, why did you select the documented design?
  (LO\_Explores)
  \item (After you have implemented another team's module, which means this
  isn't filled in until after the original deadline). What did you learn by
  implementing another team's module? Were all the details you needed in the
  documentation, or did you need to make assumptions, or ask the other team
  questions?  If your team also had another team implement one of your modules,
  what was this experience like?  Are there things in your documentation you
  could have changed to make the process go more smoothly for when an
  ``outsider'' completes some of the implementation?

  Your reflection should be about concepts learned in your SE education, such as modular decomposition, cohesion, coupling, information hiding, etc.  You should critique the quality and utility of the specification you received.

  You should create a PR for your implemented module.  A link to the PR should be part of your reflection.

\end{enumerate}

\end{document}